\documentclass[12pt]{article}
\usepackage[utf8]{inputenc}
\usepackage[T1]{fontenc}
\usepackage{hyperref}
\title{Tokano Language Resources for K-12 Curriculum Development}
\author{Compiled by Claude Sonnet 5.5 (claude-sonnet-5-5) for Andr\'as Kornai}
\date{}
\begin{document}
\maketitle

\section{Introduction}
Tokano (ISO 639-3: zuh; Glottocode toka1244) is a Trans-New Guinea language of the Lower Asaro area, Goroka District, Eastern Highlands Province, Papua New Guinea. Wikipedia gives about 6,000 speakers, citing a 1982 figure, so the number is old. Published materials are very limited.

\section{Primary Educational Resources}
No primers, textbooks or school programmes for Tokano were found in search results.

\section{Linguistic and Reference Materials}
\begin{itemize}
\item Wikipedia article on the Tokano language. It mentions a folk-tale collection translated by John Guhise (SIL, 1977): \url{https://en.wikipedia.org/wiki/Tokano_language}
\item Wikitongues page for zuh: \url{https://wikitongues.org/languages/zuh/}
\end{itemize}

\section{Religious and Community Resources}
\begin{itemize}
\item Joshua Project profile of the Tokano (Lower Asaro) people, which notes the need for scripture and teaching in Tokano: \url{https://joshuaproject.net/people_groups/15534}
\end{itemize}
Wikipedia states that portions of the Bible exist in Tokano. No direct URL for them was seen in the search results.

\section{Digital and Audio Resources}
No audio, video, corpus or software resources were found. A search result from the Max Planck Institute (DoBeS) was not examined, so nothing is claimed about it.

\section{Curriculum Development Context}
Tokano appears to be a stable community language with almost no published educational material. A curriculum would have to begin with orthography work and the collection of oral texts. It would also need coordination with local education authorities in Eastern Highlands Province.
\end{document}